%% LaRT v2 User Manual
\documentclass[english,a4paper]{kiseon_memo}
\synctex=1
\usepackage{listings}
\usepackage{longtable}
\usepackage{booktabs}
\usepackage{array}
\usepackage{babel}

\newcommand{\fdg}{\hbox{$.\!\!^\circ$}}
\newcommand{\parm}[1]{\texttt{par\%#1}}
\newcommand{\ion}[2]{#1\,{\scriptsize #2}}
\newcolumntype{L}[1]{>{\raggedright\arraybackslash}p{#1}}

\lstset{
  basicstyle=\ttfamily\small,
  commentstyle=\color{gray},
  stringstyle=\color{olive},
  frame=single,
  breaklines=true,
  numbers=left,
  numberstyle=\tiny\color{gray},
  language=[95]Fortran
}

\title{LaRT v2 User Manual}
\author{Kwang-Il Seon}
\date{2026-04-25}

\begin{document}
\maketitle
\tableofcontents
\bigskip

%=============================================================================
\section{Introduction}
%=============================================================================

\textbf{LaRT} (\textbf{La}yman-alpha \textbf{R}adiative \textbf{T}ransfer) is a Monte Carlo
radiative transfer code for computing the propagation of resonance-line photons through
astrophysical neutral-gas media.  The default transition is Lyman-$\alpha$ (Ly$\alpha$,
\ion{H}{I} 1216\,\AA), but many other UV/optical/IR lines are supported
(see \S\ref{sec:lineid}).  Dust scattering and absorption, full Stokes polarization
tracking, Hubble-flow and other bulk velocity fields, and spatially varying density,
temperature, and velocity from external data files are all supported.

Starting from v2.00, the Cartesian grid is supplemented by:
\begin{itemize}
  \item \textbf{AMR mode} — octree adaptive mesh refinement, allowing spatially varying
        resolution (\S\ref{sec:amr}).
  \item \textbf{Clumpy medium mode} — a population of compact, optically thick spherical
        clumps embedded in a transparent inter-clump medium (\S\ref{sec:clump}).
\end{itemize}

This manual covers v2.00 and v2.10.  The only user-visible change in v2.10 is the
introduction of \parm{grid\_type} as a unified selector for the three grid modes;
the old \parm{use\_amr\_grid} and \parm{use\_clump\_medium} flags remain valid for
backward compatibility.

%=============================================================================
\section{Installation and Build}
%=============================================================================

\subsection{Prerequisites}

\begin{itemize}
  \item Fortran 2003+ compiler: Intel \texttt{ifort}/\texttt{ifx} or GNU
        \texttt{gfortran}~$\geq$~7.
  \item CFITSIO library.
  \item MPI library (MPICH or OpenMPI).
\end{itemize}

\subsection{Compilation}

\begin{lstlisting}[language=bash,numbers=none]
cd LaRT_v2.00        # or LaRT_v2.10
make                 # produces LaRT.x  (pure MPI)
make DEBUG=1         # full debug flags + runtime checks
make F90=mpif90      # override compiler
make clean           # remove .o and .mod files
make cleanall        # also remove executables
make ramses2fits     # auxiliary RAMSES converter
\end{lstlisting}

Every \texttt{make} invocation recompiles all objects (no incremental builds).

\subsection{Running}

\begin{lstlisting}[language=bash,numbers=none]
mpirun -np 8 LaRT.x input.in
\end{lstlisting}

Output is written to \texttt{<base>.fits.gz}, where \texttt{<base>} is the input file
base name (without \texttt{.in}) unless \parm{out\_file} is set explicitly.

%=============================================================================
\section{Input File Format}
%=============================================================================

Input files use Fortran \emph{namelist} syntax.  All parameters are prefixed with
\texttt{par\%} and collected in a single \texttt{\&parameters} block.

\begin{lstlisting}
&parameters
 par%no_photons     = 1e6
 par%temperature    = 1.0e4      ! gas temperature [K]
 par%taumax         = 1.0e4
 par%nx = 101, par%ny = 101, par%nz = 101
 par%rmax           = 1.0
 par%source_geometry = 'point'
 par%spectral_type  = 'voigt'
 par%nxfreq         = 121
/
\end{lstlisting}

\begin{itemize}
  \item String values must be enclosed in single quotes.
  \item Multiple parameters may appear on one line, separated by commas.
  \item Lines beginning with \texttt{!} are comments.
  \item Unspecified parameters retain compiled defaults (listed in \S\ref{sec:refcard}).
\end{itemize}

%=============================================================================
\section{Grid Configuration}
\label{sec:grid}
%=============================================================================

\subsection{Cartesian Grid}

The simulation domain is a rectangular box centred at the origin.  By default
$x \in [-x_\text{max},\,+x_\text{max}]$ (and similarly for $y$ and $z$).
The cell size is $\Delta x = 2\,x_\text{max}/n_x$.

\begin{center}
\begin{tabular}{lll}
\toprule
Parameter & Default & Description \\
\midrule
\parm{nx}, \parm{ny}, \parm{nz} & 101 & Number of cells per axis \\
\parm{xmax}, \parm{ymax}, \parm{zmax} & 1.0 & Half-box size per axis (code units) \\
\parm{rmax} & $-999$ & Spherical boundary: density $= 0$ for $r > r_\text{max}$ \\
\parm{rmin} & $-999$ & Inner cavity: density $= 0$ for $r < r_\text{min}$ \\
\bottomrule
\end{tabular}
\end{center}

If \parm{input\_field} or \parm{dens\_file} is given, \texttt{(nx,ny,nz)} and
\texttt{(xmax,ymax,zmax)} are read from the FITS header automatically.

\subsection{Grid Geometry}
\label{sec:geometry}

\parm{geometry} selects the spatial symmetry used for accumulating and saving the
mean-intensity array $J$:

\begin{center}
\begin{tabular}{ll}
\toprule
Value & $J$ output geometry \\
\midrule
\texttt{'sphere'} (default) & 1-D radial profile $J(r)$ \\
\texttt{'cylinder'} & 2-D cylindrical $J(\rho,z)$ \\
\texttt{'rectangle'} & 3-D Cartesian cube $J(x,y,z)$ \\
\texttt{'plane\_atmosphere'} & 1-D slab $J(z)$ \\
\texttt{'spherical\_atmosphere'} & 1-D spherical shell $J(r)$ \\
\bottomrule
\end{tabular}
\end{center}

\subsection{Symmetry and Periodic Boundaries}

\begin{center}
\begin{tabular}{lll}
\toprule
Parameter & Default & Effect \\
\midrule
\parm{xyz\_symmetry} & \texttt{.false.} & Simulate only the $+x,+y,+z$ octant ($8\times$ RAM saving) \\
\parm{xy\_symmetry}  & \texttt{.false.} & Simulate only the $+x,+y$ quadrant ($4\times$ RAM saving) \\
\parm{xy\_periodic}  & \texttt{.false.} & Periodic in $x,y$: photons exiting one face re-enter the opposite \\
\bottomrule
\end{tabular}
\end{center}

Symmetry options are disabled automatically when peel-off is active.

\subsection{Distance Units}

\begin{center}
\begin{tabular}{lll}
\toprule
Parameter & Default & Description \\
\midrule
\parm{distance\_unit} & \texttt{''} & Physical unit: \texttt{'kpc'}, \texttt{'pc'}, \texttt{'au'}, or \texttt{''} (dimensionless) \\
\parm{distance2cm}    & 1.0 & Alternatively: 1 code unit in cm \\
\bottomrule
\end{tabular}
\end{center}

In dimensionless mode (\parm{distance\_unit}~$=$~\texttt{''}), the gas density is
determined entirely by \parm{taumax} and the grid geometry; no physical unit is needed.

\subsection{Grid Type Selection}

\begin{center}
\begin{tabular}{lll}
\toprule
Parameter & Default & Description \\
\midrule
\parm{grid\_type}$^{\dagger}$ & \texttt{'car'} & \textbf{v2.10+}: \texttt{'car'} (Cartesian), \texttt{'amr'}, or \texttt{'clump'} \\
\parm{use\_amr\_grid} & \texttt{.false.} & \textbf{v2.00}: enable AMR (deprecated in v2.10) \\
\parm{use\_clump\_medium} & \texttt{.false.} & \textbf{v2.00}: enable clump mode (deprecated in v2.10) \\
\bottomrule
\end{tabular}
\end{center}

%=============================================================================
\section{Source Configuration}
\label{sec:source}
%=============================================================================

\subsection{Source Geometry}
\label{sec:source_geometry}

\parm{source\_geometry} controls the spatial distribution of photon emission:

\begin{center}
\begin{tabular}{L{3.5cm}L{9.5cm}}
\toprule
Value & Description \\
\midrule
\texttt{'point'} & Point source at $(\parm{xs\_point},\,\parm{ys\_point},\,\parm{zs\_point})$; default $(0,0,0)$ \\
\texttt{'uniform'} & Uniform volume emissivity throughout the grid (or inside \parm{source\_rmax}) \\
\texttt{'uniform\_sphere'} & Uniform within a sphere of radius \parm{source\_rmax} \\
\texttt{'uniform\_xy'} & Uniform in the $z=0$ plane \\
\texttt{'gaussian'} & Gaussian: $\exp(-r^2/2\sigma_r^2)\exp(-z^2/2\sigma_z^2)$; widths \parm{source\_rscale}, \parm{source\_zscale} \\
\texttt{'exponential'} & Exponential disc: $\exp(-r/r_s)\exp(-|z|/z_s)$; scales \parm{source\_rscale}, \parm{source\_zscale} \\
\texttt{'exponential\_sphere'} & Exponential clipped to sphere of radius \parm{source\_rmax} \\
\texttt{'exponential\_cylinder'} & Exponential clipped to a cylinder \\
\texttt{'sersic'} / \texttt{'ssh'} & S\'{e}rsic profile; index \parm{sersic\_m}, half-light radius \parm{Reff} \\
\texttt{'diffuse\_emissivity'} & 3-D emissivity read from \parm{emiss\_file} \\
\texttt{'star\_file'} & Individual point sources from text file \parm{star\_file} (columns: $x,y,z,L$) \\
\bottomrule
\end{tabular}
\end{center}

Associated parameters:

\begin{center}
\begin{tabular}{lll}
\toprule
Parameter & Default & Description \\
\midrule
\parm{xs\_point}, \parm{ys\_point}, \parm{zs\_point} & 0.0 & Point source position (code units) \\
\parm{source\_rmax}   & $-999$ & Outer truncation radius for extended sources \\
\parm{source\_rscale} & $-999$ & Radial scale length \\
\parm{source\_zscale} & 0.0   & Vertical scale height \\
\parm{sersic\_m}      & 4.0   & S\'{e}rsic index \\
\parm{Reff}           & 1.0   & Effective (half-light) radius \\
\parm{comoving\_source} & \texttt{.true.} & Source photons co-moving with the local gas \\
\bottomrule
\end{tabular}
\end{center}

\subsection{Spectral Type}
\label{sec:spectral_type}

\parm{spectral\_type} sets the frequency distribution of emitted photons:

\begin{center}
\begin{tabular}{L{3.5cm}L{9.5cm}}
\toprule
Value & Description \\
\midrule
\texttt{'monochromatic'} & Delta function at line centre ($x = 0$) \\
\texttt{'voigt'} & Voigt profile at the local gas temperature \\
\texttt{'voigt0'} & Voigt profile at the fixed temperature \parm{temperature0} \\
\texttt{'gaussian'} & Gaussian with width \parm{gaussian\_width\_vel} [km\,s$^{-1}$] \\
\texttt{'continuum'} & Flat spectrum over [\parm{xfreq\_min}, \parm{xfreq\_max}] \\
\texttt{'line\_prof\_file'} & Arbitrary profile read from \parm{line\_prof\_file} \\
\bottomrule
\end{tabular}
\end{center}

\subsection{Line Identification}
\label{sec:lineid}

The atomic transition is selected with \parm{line\_id} (default: Ly$\alpha$).

\begin{center}
\begin{tabular}{lll}
\toprule
\parm{line\_id} & Transition & Wavelength \\
\midrule
(default) & \ion{H}{I} Ly$\alpha$ & 1216\,\AA \\
\texttt{'HeI\_10833'} & He~I & 10833\,\AA \\
\texttt{'CIV\_1548'} & \ion{C}{IV} & 1548\,\AA \\
\texttt{'NV\_1239'} & \ion{N}{V} & 1239\,\AA \\
\texttt{'OVI\_1032'} & \ion{O}{VI} & 1032\,\AA \\
\texttt{'MgII\_2796'} & \ion{Mg}{II} & 2796\,\AA \\
\texttt{'SiII\_1527'} & \ion{Si}{II} & 1527\,\AA \\
\texttt{'SiII\_1260'} & \ion{Si}{II} & 1260\,\AA \\
\texttt{'SiII\_1193'} / \texttt{'SiII\_1190'} & \ion{Si}{II} & 1193/1190\,\AA \\
\texttt{'SiII\_1304'} & \ion{Si}{II} & 1304\,\AA \\
\texttt{'FeII\_UV1'} / \texttt{'FeII\_2600'} & Fe~II & 2600\,\AA \\
\texttt{'FeII\_UV2'} / \texttt{'FeII\_2383'} & Fe~II & 2383\,\AA \\
\texttt{'FeII\_UV3'} / \texttt{'FeII\_2344'} & Fe~II & 2344\,\AA \\
\texttt{'FeII\_2250'} & Fe~II & 2250\,\AA \\
\texttt{'FeII\_2261'} & Fe~II & 2261\,\AA \\
\bottomrule
\end{tabular}
\end{center}

\subsection{Frequency and Wavelength Grid}
\label{sec:freqgrid}

Internally, photon frequencies are tracked in the dimensionless Doppler variable:
\begin{equation}
  x = \frac{\nu - \nu_0}{\Delta\nu_D}, \qquad
  \Delta\nu_D = \frac{\nu_0}{c}\sqrt{\frac{2kT}{m}},
\end{equation}
where $\nu_0$ is the line-centre frequency, $T$ is the local gas temperature, and $m$
is the atomic/ionic mass.  Output spectra are binned onto a uniform grid in $x$ by
default.  A wavelength or velocity grid can be requested instead:

\begin{center}
\begin{tabular}{lll}
\toprule
Parameter & Default & Description \\
\midrule
\parm{nxfreq}   & 121  & Number of spectral bins \\
\parm{xfreq\_min}, \parm{xfreq\_max} & auto & Limits of $x$-grid \\
\parm{xfreq0}   & 0.0  & Frequency offset applied to all emitted photons \\
\parm{nwavelength} & 0 & Number of wavelength bins (overrides $x$-grid) \\
\parm{wavelength\_min}, \parm{wavelength\_max} & --- & [\AA] \\
\parm{nvelocity} & 0   & Number of velocity bins (overrides $x$-grid) \\
\parm{velocity\_min}, \parm{velocity\_max} & --- & [km\,s$^{-1}$] \\
\bottomrule
\end{tabular}
\end{center}

The range is set automatically from the mean medium temperature when no limits are provided.

%=============================================================================
\section{Medium Properties}
\label{sec:medium}
%=============================================================================

\subsection{Optical Depth Normalization}

The overall opacity is scaled by one of the following (specify at most one):

\begin{center}
\begin{tabular}{L{3.2cm}L{10cm}}
\toprule
Parameter & Description \\
\midrule
\parm{taumax}   & Line-centre $\tau_0$ from the grid centre to the $+z$ boundary along the $z$-axis \\
\parm{tauhomo}  & Same, assuming a uniform density field (scales an inhomogeneous field proportionally) \\
\parm{N\_HImax} & \ion{H}{I} column density from centre to $+z$ boundary [cm$^{-2}$] \\
\parm{N\_HIhomo} & Same, assuming uniform density \\
\bottomrule
\end{tabular}
\end{center}

For AMR grids, \parm{taumax} is the pole-to-face optical depth (box centre $\to$ $+z$ face
along $z$), consistent with the Cartesian convention for a sphere of radius $R = L/2$.

\subsection{Temperature}

\begin{center}
\begin{tabular}{lll}
\toprule
Parameter & Default & Description \\
\midrule
\parm{temperature}  & $10^4$\,K & Uniform gas temperature \\
\parm{temperature0} & $= T$ & Temperature for \texttt{spectral\_type='voigt0'} source \\
\parm{bturb}        & 0.0   & Turbulent velocity dispersion added in quadrature to thermal width [km\,s$^{-1}$] \\
\bottomrule
\end{tabular}
\end{center}

A spatially varying temperature field can be loaded from a FITS file via \parm{temp\_file}
(see the external data table in \S\ref{sec:density}).

\subsection{Density}
\label{sec:density}

\textbf{Analytic profiles.}
Without external files, the density is set analytically via the opacity scale parameters
(\parm{taumax} etc.) combined with an optional radial profile:

\begin{center}
\begin{tabular}{lll}
\toprule
Profile type & Key parameter & Functional form \\
\midrule
Uniform (sphere or box) & \parm{rmax} & Constant $n_\mathrm{H}$ inside $r < r_\text{max}$ \\
Spherical shell & \parm{rmin}, \parm{rmax} & Constant $n_\mathrm{H}$ between $r_\text{min}$ and $r_\text{max}$ \\
Exponential & \parm{density\_rscale} & $n \propto \exp(-r/r_s)$ \\
\bottomrule
\end{tabular}
\end{center}

\textbf{External FITS data.}

\begin{center}
\begin{tabular}{L{3.2cm}L{10cm}}
\toprule
Parameter & Description \\
\midrule
\parm{dens\_file}  & 3-D \ion{H}{I} number density cube [cm$^{-3}$] \\
\parm{temp\_file}  & 3-D temperature cube [K] \\
\parm{velo\_file}  & 4-D velocity cube $(v_x,v_y,v_z)$ [km\,s$^{-1}$] \\
\parm{emiss\_file} & 3-D emissivity cube [photons\,s$^{-1}$\,cm$^{-3}$] \\
\parm{input\_field} & Base name: reads \texttt{<base>.dens.fits.gz}, \texttt{.temp.}, \texttt{.velo.} \\
\bottomrule
\end{tabular}
\end{center}

If \parm{taumax} (or \parm{N\_HImax}) is specified alongside a density file, the density
is rescaled so that the target optical depth is achieved.

\subsection{Velocity Fields}
\label{sec:velocity}

\parm{velocity\_type} selects the bulk velocity model.  The default (empty string) is a
static medium.

\begin{center}
\begin{tabular}{L{3.8cm}L{9.5cm}}
\toprule
Value & Velocity law \\
\midrule
\texttt{'hubble'} & $\mathbf{v}(r) = V_\text{exp}\,(r/r_\text{max})\,\hat{r}$;\quad $V_\text{exp}>0$: outflow, $<0$: infall \\
\texttt{'constant\_radial'} & $\mathbf{v}(r) = V_\text{exp}\,\hat{r}$ (constant magnitude radial flow) \\
\texttt{'parallel\_velocity'} & Uniform $\mathbf{v} = (V_x,V_y,V_z)$ throughout the box \\
\texttt{'ssh'} & Bilinear radial profile (Song, Seon \& Hwang 2020); see below \\
\texttt{'rotating\_solid\_body'} & Pure solid-body rotation in the $xy$-plane: $v_x = -V_\text{rot}\,y/r_\text{max}$, $v_y = V_\text{rot}\,x/r_\text{max}$; speed $\propto \rho/r_\text{max}$, reaching $V_\text{rot}$ at $\rho = r_\text{max}$ (Garavito-Camargo et al.\ 2014) \\
\texttt{'rotating\_galaxy\_halo'} & Solid-body rise to $V_\text{rot}$ for $\rho < r_\text{inner}$, then flat at $V_\text{rot}$ for $\rho \geq r_\text{inner}$; rotation in the $xy$-plane only (Kim et al.) \\
\bottomrule
\end{tabular}
\end{center}

The \texttt{'ssh'} profile rises linearly to $V_\text{peak}$ at $r_\text{peak}$, then
shifts by $\Delta V$ toward $r_\text{max}$:
\begin{equation}
  v(r) = \begin{cases}
    V_\text{peak}\,r/r_\text{peak} & r < r_\text{peak} \\
    \left[V_\text{peak} + \Delta V\,\frac{r - r_\text{peak}}{r_\text{max} - r_\text{peak}}\right]\hat{r} & r \geq r_\text{peak}.
  \end{cases}
\end{equation}

A velocity FITS file (\parm{velo\_file}) may be combined with a systematic \parm{velocity\_type}.

\begin{center}
\begin{tabular}{lll}
\toprule
Parameter & Default & Description \\
\midrule
\parm{Vexp}       & 0.0  & Expansion/infall speed for \texttt{hubble} and \texttt{constant\_radial} [km\,s$^{-1}$] \\
\parm{Vx}, \parm{Vy}, \parm{Vz} & 0.0 & Uniform velocity components for \texttt{parallel\_velocity} [km\,s$^{-1}$] \\
\parm{Vpeak}      & 0.0  & Peak velocity for \texttt{ssh} [km\,s$^{-1}$] \\
\parm{rpeak}      & 0.0  & Radius at peak for \texttt{ssh} (code units) \\
\parm{DeltaV}     & 0.0  & Velocity increment beyond $r_\text{peak}$ for \texttt{ssh} [km\,s$^{-1}$] \\
\parm{Vrot}       & 0.0  & Maximum rotation speed for \texttt{rotating\_solid\_body} and \texttt{rotating\_galaxy\_halo} [km\,s$^{-1}$] \\
\parm{rinner}     & 0.0  & Transition radius for \texttt{rotating\_galaxy\_halo}: solid-body rise for $\rho < r_\text{inner}$, flat at $V_\text{rot}$ beyond \\
\parm{Omega}      & 28.0 & Shear rate [(km\,s$^{-1}$)\,kpc$^{-1}$] for TIGRESS data \\
\parm{q}          & 0.0  & Shear factor multiplier for TIGRESS data (effective $\Omega \leftarrow q\,\Omega\,L_x$) \\
\parm{recoil}     & \texttt{.false.} & Include photon-recoil frequency shift in scattering \\
\parm{core\_skip} & \texttt{.false.} & Enable core-skipping acceleration algorithm \\
\bottomrule
\end{tabular}
\end{center}

%=============================================================================
\section{AMR Grid Mode}
\label{sec:amr}
%=============================================================================

\subsection{Enabling AMR Mode}

\textbf{v2.00:}
\begin{lstlisting}
par%use_amr_grid = .true.
par%amr_type     = 'generic'   ! always specify explicitly
par%amr_file     = 'grid.dat'
\end{lstlisting}

\textbf{v2.10 (preferred):}
\begin{lstlisting}
par%grid_type = 'amr'
par%amr_type  = 'generic'
par%amr_file  = 'grid.dat'
\end{lstlisting}

\begin{center}
\begin{tabular}{lll}
\toprule
Parameter & Default & Description \\
\midrule
\parm{amr\_type}    & \texttt{'ramses'} & Format: \texttt{'generic'} (text) or \texttt{'ramses'} (binary snapshot) \\
\parm{amr\_file}    & \texttt{''} & Path to AMR data file or RAMSES snapshot directory \\
\parm{amr\_snapnum} & 1 & Snapshot number for RAMSES output \\
\bottomrule
\end{tabular}
\end{center}

\textbf{Note:} the compiled default for \parm{amr\_type} is \texttt{'ramses'}.  Always
specify \texttt{par\%amr\_type = 'generic'} explicitly when using the text format.

\subsection{Generic AMR Data Format}

Plain text: a header followed by one line per leaf cell.

\begin{lstlisting}[language={}]
# comment lines start with #
BOXLEN  100.0            ! box side length in code units
NLEAF   8192             ! number of leaf cells
! columns: x  y  z  level  nH[cm^-3]  T[K]  vx[km/s]  vy[km/s]  vz[km/s]
50.5  50.5  50.5  1  1.0e-3  1.0e4  0.0  0.0  0.0
...
\end{lstlisting}

Leaf-cell coordinates are given with the box origin at the corner, $[0,\,\text{BOXLEN}]$.
The reader internally re-centres to $[-L/2,\,+L/2]$; \emph{do not} pre-shift the data.

\subsection{Tau Normalization Convention}

For AMR, \parm{taumax} is the line-centre optical depth from the box centre to the
$+z$ face along the $z$-axis.  This matches the Cartesian sphere convention where
\parm{taumax} is the centre-to-surface optical depth, provided $r_\text{max} = L/2$.

\subsection{Source Geometries in AMR Mode}

In v2.00, AMR supports \texttt{'point'} and \texttt{'uniform'}.  In v2.10, most
Cartesian geometries (Gaussian, exponential, S\'{e}rsic, \texttt{star\_file}, etc.)
are also available; refer to \S\ref{sec:source_geometry} and the v2.10 release notes.

%=============================================================================
\section{Clumpy Medium Mode}
\label{sec:clump}
%=============================================================================

\subsection{Overview}

Clumpy medium mode places $N$ non-overlapping spherical clumps of radius $r_c$ inside a
spherical outer domain of radius $R$.  The inter-clump medium is vacuum.  All opacity
resides inside the clumps.  Clump positions are drawn by Random Sequential Adsorption (RSA).

\subsection{Enabling Clump Mode}

\textbf{v2.00:}
\begin{lstlisting}
par%use_clump_medium = .true.
\end{lstlisting}

\textbf{v2.10 (preferred):}
\begin{lstlisting}
par%grid_type = 'clump'
\end{lstlisting}

\subsection{Clump Parameters}

\begin{center}
\begin{tabular}{lll}
\toprule
Parameter & Default & Description \\
\midrule
\parm{rmax}            & 1.0  & Outer sphere radius (code units) \\
\parm{clump\_radius}   & ---  & Individual clump radius (same units) \\
\parm{clump\_tau0}     & ---  & Line-centre $\tau_0$ from clump centre to its surface \\
\parm{clump\_f\_cov}   & ---  & Covering factor $C_f = N r_c^2 / R^2$ \\
\parm{clump\_f\_vol}   & ---  & Volume filling factor $f_V = N(r_c/R)^3$ \\
\parm{clump\_N\_clumps} & --- & Total number of clumps $N$ \\
\parm{temperature}     & $10^4$\,K & Clump gas temperature \\
\parm{clump\_sigma\_v} & 0.0  & Gaussian $\sigma$ of random bulk clump velocities [km\,s$^{-1}$] \\
\parm{save\_clump\_info} & \texttt{.false.} & Write clump positions and velocities to a FITS binary table \\
\bottomrule
\end{tabular}
\end{center}

Specify exactly one of \parm{clump\_f\_cov}, \parm{clump\_f\_vol}, or \parm{clump\_N\_clumps}.

\subsection{Clump Velocity}

A systematic flow can be added on top of the random component (\parm{clump\_sigma\_v})
via \parm{velocity\_type}.  Supported values in clump mode: \texttt{'hubble'},
\texttt{'constant\_radial'}, \texttt{'parallel\_velocity'}, \texttt{'ssh'},
\texttt{'rotating\_solid\_body'}, \texttt{'rotating\_galaxy\_halo'}.

%=============================================================================
\section{Observers and Peel-off}
\label{sec:observer}
%=============================================================================

The peel-off (peeling-off) technique shoots a photon toward the observer at every
scattering event, allowing spectral images to be constructed without noise.  It is
activated by setting \parm{nxim} and \parm{nyim} to non-zero values.

\subsection{Observer Placement}

\textbf{Direction vector:}
\begin{lstlisting}
par%obsx = 0.0, par%obsy = 0.0, par%obsz = 1.0  ! look down +z axis
par%distance = 100.0                               ! observer-centre separation
\end{lstlisting}

\textbf{Spherical angles:}
\begin{lstlisting}
par%alpha = 45.0   ! azimuthal angle [deg]
par%beta  = 30.0   ! elevation angle [deg]
par%distance = 100.0
\end{lstlisting}

\parm{gamma} sets the rotation of the detector plane about the line of sight.
If \parm{distance} is omitted it defaults to $100 \times r_\text{max}$ (or the maximum
box half-extent).

\subsection{Multiple Observers}

Multiple observers are specified as arrays:
\begin{lstlisting}
par%alpha = 0.0  45.0  90.0  135.0
par%beta  = 0.0   0.0   0.0    0.0
\end{lstlisting}
Up to 99 simultaneous observers are supported.

\subsection{Image Parameters}

\begin{center}
\begin{tabular}{lll}
\toprule
Parameter & Default & Description \\
\midrule
\parm{nxim}, \parm{nyim} & 0    & Image pixel counts; peel-off enabled when $> 0$ \\
\parm{dxim}, \parm{dyim} & auto & Pixel angular size [deg] (defaults to cover the full grid) \\
\bottomrule
\end{tabular}
\end{center}

When \parm{dxim}/\parm{dyim} are not set, the pixel size is chosen to cover the full
sphere:
\begin{equation}
  \delta\theta = \frac{1}{N_x/2}\arcsin\!\left(\frac{r_\text{max}}{d}\right).
\end{equation}

\subsection{Peel-off Output Flags}

\begin{center}
\begin{tabular}{lll}
\toprule
Parameter & Default & Description \\
\midrule
\parm{save\_peeloff}    & \texttt{.false.} & Save 1-D peel-off spectra per observer \\
\parm{save\_peeloff\_2D} & \texttt{.false.} & Save peel-off spectral image cubes ($x$-$y$-$\nu$) \\
\parm{save\_peeloff\_3D} & \texttt{.false.} & Save full 3-D peel-off data \\
\parm{save\_direc0}     & \texttt{.false.} & Save the peel-off image assuming a fully transparent medium (no $e^{-\tau}$ attenuation to the observer); provides the unattenuated direct contribution as a reference \\
\bottomrule
\end{tabular}
\end{center}

%=============================================================================
\section{Output Files}
\label{sec:output}
%=============================================================================

All outputs are gzip-compressed FITS files (\texttt{.fits.gz}).

\subsection{Primary Output Arrays}

\begin{center}
\begin{tabular}{lll}
\toprule
Array & FITS extension & Description \\
\midrule
\texttt{Jout} & \texttt{JOUT} & Emergent spectrum integrated over all escape directions \\
\texttt{Jin}  & \texttt{JIN}  & Injected (intrinsic) source spectrum \\
\texttt{Jabs} & \texttt{JABS} & Spectrum of photons absorbed by dust \\
\bottomrule
\end{tabular}
\end{center}

\textbf{Jout} is the main science output.  It is normalized per unit frequency interval and
per unit area of the bounding sphere.  In dimensionless units (no \parm{distance\_unit}):
\begin{equation}
  J^\mathrm{out}(x) = \frac{\text{escaped photon weight in bin } x}
                            {n_\mathrm{ph}\;\Delta x\;2\pi\;\pi R^2},
\end{equation}
where $R = r_\text{max}$ (Cartesian) or $R = L_\text{box}/2$ (AMR), and
$\Delta x = 1$ in units of the Doppler width.

\textbf{Jin} requires \parm{save\_Jin}~$=$~\texttt{.true.}

\textbf{Jabs} requires \parm{save\_Jabs}~$=$~\texttt{.true.}

The total flux is conserved: $\int J^\mathrm{out}\,dx + \int J^\mathrm{abs}\,dx = \int J^\mathrm{in}\,dx$
(in the absence of other losses).

\subsection{Output Control Parameters}

\begin{center}
\begin{tabular}{lll}
\toprule
Parameter & Default & Description \\
\midrule
\parm{out\_file}   & \texttt{''} & Output file name (defaults to \texttt{<input>.fits.gz}) \\
\parm{out\_merge}  & \texttt{.false.} & Accumulate into existing output instead of overwriting \\
\parm{out\_bitpix} & $-32$ & FITS precision: $-32$\,=\,float32, $-64$\,=\,float64 \\
\parm{save\_Jin}   & \texttt{.false.} & Save injected spectrum \\
\parm{save\_Jabs}  & \texttt{.false.} & Save dust-absorbed spectrum \\
\parm{save\_all}   & \texttt{.false.} & Save full 3-D $J(\nu,x,y,z)$ array \\
\parm{save\_radial\_profile} & \texttt{.false.} & Save 1-D radial $J$ profile \\
\parm{save\_sightline\_tau} & \texttt{.false.} & Save optical depth along sightlines \\
\parm{intensity\_unit} & \texttt{''} & Physical unit label written to the FITS header \\
\bottomrule
\end{tabular}
\end{center}

\parm{out\_merge}~$=$~\texttt{.true.} is useful for noise reduction: run the same setup
multiple times and merge the results.  All input parameters must be identical across runs.

%=============================================================================
\section{Dust and Polarization}
\label{sec:dust}
%=============================================================================

\begin{center}
\begin{tabular}{lll}
\toprule
Parameter & Default & Description \\
\midrule
\parm{DGR}    & 0.0     & Dust-to-gas ratio relative to the Milky Way value \\
\parm{albedo} & 0.32536 & Single-scattering albedo $\omega$ \\
\parm{hgg}    & 0.67617 & Henyey--Greenstein asymmetry parameter $g$ \\
\parm{use\_stokes}      & \texttt{.false.} & Enable full Stokes polarization RT \\
\parm{scatt\_mat\_file} & \texttt{''} & Mueller matrix file for polarized dust scattering \\
\parm{use\_reduced\_wgt} & \texttt{.false.} & Reduce photon weight by albedo instead of destroying absorbed photons \\
\bottomrule
\end{tabular}
\end{center}

When \parm{DGR}~$= 0$ (default), there is no dust and all photon weight is conserved.
Non-zero \parm{DGR} adds dust extinction and Henyey--Greenstein scattering.  When
\parm{use\_stokes}~$=$~\texttt{.true.}, a Mueller matrix file must be provided;
files for the Weingartner--Draine Milky Way model are included in the \texttt{data/}
directory.

%=============================================================================
\section{Parallelism}
\label{sec:parallel}
%=============================================================================

\begin{center}
\begin{tabular}{lll}
\toprule
Parameter & Default & Description \\
\midrule
\parm{no\_photons}        & $10^6$ & Total photon packets (floating-point, converted to int64) \\
\parm{use\_master\_slave} & \texttt{.true.} & Dynamic load balancing via master--worker algorithm \\
\parm{num\_send\_at\_once} & 100  & Photons dispatched to each worker per batch \\
\parm{nprint}  & $10^5$ & Progress-report interval [photons] \\
\parm{iseed}   & auto   & Random seed; defaults derived from MPI rank and system clock \\
\bottomrule
\end{tabular}
\end{center}

With \parm{use\_master\_slave}~$=$~\texttt{.true.}, one MPI rank acts as dispatcher and
the rest as workers.  This balances workloads when scattering rates vary across photon
packets.  With \texttt{.false.}, each rank processes an equal share of \parm{no\_photons}.

%=============================================================================
\section{Complete Parameter Reference}
\label{sec:refcard}
%=============================================================================

Table~\ref{tab:refcard} lists commonly used parameters.  Parameters marked
$\dagger$ are new or changed in v2.10.

\begin{longtable}{L{3.6cm} L{2.4cm} L{7.4cm}}
\caption{LaRT parameter reference.\label{tab:refcard}} \\
\toprule
\textbf{Parameter} & \textbf{Default} & \textbf{Description} \\
\midrule
\endfirsthead
\multicolumn{3}{l}{\small\textit{(continued)}} \\[2pt]
\toprule
\textbf{Parameter} & \textbf{Default} & \textbf{Description} \\
\midrule
\endhead
\midrule
\multicolumn{3}{r}{\small\textit{continued on next page}} \\
\endfoot
\bottomrule
\endlastfoot

\multicolumn{3}{l}{\textit{Simulation control}} \\[2pt]
\texttt{par\%no\_photons}  & $10^6$ & Number of photon packets \\
\texttt{par\%nprint}       & $10^5$ & Progress interval [photons] \\
\texttt{par\%iseed}        & auto   & Random seed \\[6pt]

\multicolumn{3}{l}{\textit{Cartesian grid}} \\[2pt]
\texttt{par\%nx}, \texttt{par\%ny}, \texttt{par\%nz} & 101 & Cells per axis \\
\texttt{par\%xmax}, \texttt{par\%ymax}, \texttt{par\%zmax} & 1.0 & Half-box size (code units) \\
\texttt{par\%rmax}  & $-999$ & Spherical boundary radius \\
\texttt{par\%rmin}  & $-999$ & Inner cavity radius \\
\texttt{par\%geometry} & \texttt{'sphere'} & $J$ output geometry (\S\ref{sec:geometry}) \\
\texttt{par\%xyz\_symmetry} & \texttt{.false.} & $+x,+y,+z$ octant only \\
\texttt{par\%xy\_symmetry}  & \texttt{.false.} & $+x,+y$ quadrant only \\
\texttt{par\%xy\_periodic}  & \texttt{.false.} & Periodic in $x,y$ \\
\texttt{par\%distance\_unit} & \texttt{''} & Physical unit (\texttt{'kpc'}, \texttt{'pc'}, \texttt{'au'}) \\
\texttt{par\%distance2cm}   & 1.0 & 1 code unit in cm \\[6pt]

\multicolumn{3}{l}{\textit{Grid mode}} \\[2pt]
\texttt{par\%grid\_type}$^\dagger$ & \texttt{'car'} & v2.10+: \texttt{'car'}, \texttt{'amr'}, \texttt{'clump'} \\
\texttt{par\%use\_amr\_grid} & \texttt{.false.} & v2.00 AMR flag \\
\texttt{par\%use\_clump\_medium} & \texttt{.false.} & v2.00 clump flag \\[6pt]

\multicolumn{3}{l}{\textit{Source}} \\[2pt]
\texttt{par\%source\_geometry} & \texttt{'point'} & See \S\ref{sec:source_geometry} \\
\texttt{par\%xs\_point}, \texttt{par\%ys\_point}, \texttt{par\%zs\_point} & 0.0 & Point source position \\
\texttt{par\%source\_rmax}    & $-999$ & Source truncation radius \\
\texttt{par\%source\_rscale}  & $-999$ & Radial scale length \\
\texttt{par\%source\_zscale}  & 0.0   & Vertical scale height \\
\texttt{par\%sersic\_m}       & 4.0   & S\'{e}rsic index \\
\texttt{par\%Reff}            & 1.0   & Effective radius \\
\texttt{par\%comoving\_source} & \texttt{.true.} & Source co-moving with local gas \\
\texttt{par\%emiss\_file}     & \texttt{''} & Emissivity FITS cube \\
\texttt{par\%star\_file}      & \texttt{''} & Point-star list $(x,y,z,L)$ \\[6pt]

\multicolumn{3}{l}{\textit{Line and spectrum}} \\[2pt]
\texttt{par\%line\_id}        & Ly$\alpha$ & Atomic transition (\S\ref{sec:lineid}) \\
\texttt{par\%spectral\_type}  & \texttt{'monochromatic'} & Emission frequency distribution \\
\texttt{par\%temperature}     & $10^4$\,K & Gas temperature \\
\texttt{par\%temperature0}    & $= T$    & Source temperature for \texttt{voigt0} \\
\texttt{par\%bturb}           & 0.0  & Turbulent dispersion [km\,s$^{-1}$] \\
\texttt{par\%gaussian\_width\_vel} & 0.0 & Gaussian line width [km\,s$^{-1}$] \\
\texttt{par\%xfreq0}          & 0.0  & Frequency offset for emitted photons \\
\texttt{par\%line\_prof\_file} & \texttt{''} & Custom line profile file \\
\texttt{par\%line\_prof\_file\_type} & 0 & 0: $(x,I)$; 1: $(\lambda,I)$ \\[6pt]

\multicolumn{3}{l}{\textit{Frequency grid}} \\[2pt]
\texttt{par\%nxfreq}       & 121  & Spectral bins \\
\texttt{par\%xfreq\_min}, \texttt{par\%xfreq\_max} & auto & $x$-grid limits \\
\texttt{par\%nwavelength}  & 0    & Wavelength bins (overrides $x$-grid) \\
\texttt{par\%wavelength\_min}, \texttt{par\%wavelength\_max} & --- & [\AA] \\
\texttt{par\%nvelocity}    & 0    & Velocity bins (overrides $x$-grid) \\
\texttt{par\%velocity\_min}, \texttt{par\%velocity\_max} & --- & [km\,s$^{-1}$] \\[6pt]

\multicolumn{3}{l}{\textit{Optical depth and density}} \\[2pt]
\texttt{par\%taumax}       & ---  & Line-centre $\tau_0$: centre $\to$ $+z$ boundary \\
\texttt{par\%tauhomo}      & ---  & Same, uniform-density equivalent \\
\texttt{par\%N\_HImax}     & ---  & Column density: centre $\to$ $+z$ [cm$^{-2}$] \\
\texttt{par\%N\_HIhomo}    & ---  & Same, uniform-density equivalent \\
\texttt{par\%density\_rscale} & $-999$ & Exponential density scale radius \\[6pt]

\multicolumn{3}{l}{\textit{External data}} \\[2pt]
\texttt{par\%dens\_file}   & \texttt{''} & Density FITS cube [cm$^{-3}$] \\
\texttt{par\%temp\_file}   & \texttt{''} & Temperature FITS cube [K] \\
\texttt{par\%velo\_file}   & \texttt{''} & Velocity FITS cube [km\,s$^{-1}$] \\
\texttt{par\%input\_field} & \texttt{''} & Base name for \texttt{.dens/.temp/.velo.fits.gz} \\[6pt]

\multicolumn{3}{l}{\textit{Velocity}} \\[2pt]
\texttt{par\%velocity\_type} & \texttt{''} & Velocity law (\S\ref{sec:velocity}) \\
\texttt{par\%Vexp}  & 0.0 & Expansion/infall speed [km\,s$^{-1}$] \\
\texttt{par\%Vx}, \texttt{par\%Vy}, \texttt{par\%Vz} & 0.0 & Uniform velocity [km\,s$^{-1}$] \\
\texttt{par\%Vpeak} & 0.0 & Peak velocity for SSH [km\,s$^{-1}$] \\
\texttt{par\%rpeak} & 0.0 & Radius at peak for SSH \\
\texttt{par\%DeltaV} & 0.0 & Velocity increment for SSH [km\,s$^{-1}$] \\
\texttt{par\%Vrot}  & 0.0 & Rotation speed [km\,s$^{-1}$] \\
\texttt{par\%rinner} & 0.0 & Transition radius for \texttt{rotating\_galaxy\_halo} \\
\texttt{par\%Omega} & 28.0 & Shear rate [(km\,s$^{-1}$)\,kpc$^{-1}$] for TIGRESS data \\
\texttt{par\%q}    & 0.0 & Shear factor for TIGRESS ($\Omega \leftarrow q\,\Omega\,L_x$) \\
\texttt{par\%recoil} & \texttt{.false.} & Photon recoil in scattering \\
\texttt{par\%core\_skip} & \texttt{.false.} & Core-skipping acceleration \\[6pt]

\multicolumn{3}{l}{\textit{AMR}} \\[2pt]
\texttt{par\%amr\_type}    & \texttt{'ramses'} & \texttt{'generic'} or \texttt{'ramses'} \\
\texttt{par\%amr\_file}    & \texttt{''} & AMR data path \\
\texttt{par\%amr\_snapnum} & 1 & RAMSES snapshot number \\[6pt]

\multicolumn{3}{l}{\textit{Clumpy medium}} \\[2pt]
\texttt{par\%clump\_radius}   & --- & Clump radius \\
\texttt{par\%clump\_tau0}     & --- & Line-centre $\tau_0$ per clump \\
\texttt{par\%clump\_f\_cov}   & --- & Covering factor $C_f$ \\
\texttt{par\%clump\_f\_vol}   & --- & Volume filling factor $f_V$ \\
\texttt{par\%clump\_N\_clumps} & --- & Number of clumps $N$ \\
\texttt{par\%clump\_sigma\_v} & 0.0 & Random velocity dispersion [km\,s$^{-1}$] \\
\texttt{par\%save\_clump\_info} & \texttt{.false.} & Write clump table to FITS \\[6pt]

\multicolumn{3}{l}{\textit{Observers and peel-off}} \\[2pt]
\texttt{par\%nxim}, \texttt{par\%nyim} & 0 & Image pixels; peel-off active when $> 0$ \\
\texttt{par\%dxim}, \texttt{par\%dyim} & auto & Pixel size [deg] \\
\texttt{par\%distance}  & $100\,r_\text{max}$ & Observer--centre distance \\
\texttt{par\%obsx}, \texttt{par\%obsy}, \texttt{par\%obsz} & $(0,0,1)$ & Observer direction \\
\texttt{par\%alpha}, \texttt{par\%beta} & --- & Azimuth, elevation [deg] \\
\texttt{par\%gamma}     & 0.0 & Detector-plane rotation [deg] \\
\texttt{par\%save\_peeloff}    & \texttt{.false.} & 1-D peel-off spectra \\
\texttt{par\%save\_peeloff\_2D} & \texttt{.false.} & Peel-off spectral images \\
\texttt{par\%save\_peeloff\_3D} & \texttt{.false.} & Full 3-D peel-off cube \\
\texttt{par\%save\_direc0}    & \texttt{.false.} & Peel-off image without $e^{-\tau}$ (transparent-medium reference) \\[6pt]

\multicolumn{3}{l}{\textit{Output}} \\[2pt]
\texttt{par\%out\_file}    & \texttt{''} & Output FITS file name \\
\texttt{par\%out\_merge}   & \texttt{.false.} & Accumulate into existing file \\
\texttt{par\%out\_bitpix}  & $-32$ & FITS precision \\
\texttt{par\%save\_Jin}    & \texttt{.false.} & Save injected spectrum \\
\texttt{par\%save\_Jabs}   & \texttt{.false.} & Save absorbed spectrum \\
\texttt{par\%save\_all}    & \texttt{.false.} & Save full 3-D $J(\nu,x,y,z)$ \\
\texttt{par\%save\_radial\_profile} & \texttt{.false.} & Save 1-D radial $J$ \\
\texttt{par\%save\_sightline\_tau}  & \texttt{.false.} & Save sightline $\tau$ \\
\texttt{par\%intensity\_unit} & \texttt{''} & FITS header unit label \\[6pt]

\multicolumn{3}{l}{\textit{Dust}} \\[2pt]
\texttt{par\%DGR}    & 0.0     & Dust-to-gas ratio \\
\texttt{par\%albedo} & 0.32536 & Single-scattering albedo \\
\texttt{par\%hgg}    & 0.67617 & Henyey--Greenstein asymmetry $g$ \\
\texttt{par\%use\_stokes} & \texttt{.false.} & Polarization RT \\
\texttt{par\%scatt\_mat\_file} & \texttt{''} & Mueller matrix file \\
\texttt{par\%use\_reduced\_wgt} & \texttt{.false.} & Reduce weight instead of destroy \\[6pt]

\multicolumn{3}{l}{\textit{Parallelism}} \\[2pt]
\texttt{par\%use\_master\_slave} & \texttt{.true.} & Dynamic load balancing \\
\texttt{par\%num\_send\_at\_once} & 100 & Photons per worker batch \\
\end{longtable}

%=============================================================================
\section{Example Input Files}
\label{sec:examples}
%=============================================================================

\subsection{Uniform Sphere, Ly$\alpha$, Static}

\begin{lstlisting}
&parameters
 par%no_photons     = 1e6
 par%temperature    = 1.0e4
 par%taumax         = 1.0e4
 par%rmax           = 1.0
 par%nx = 101, par%ny = 101, par%nz = 101
 par%source_geometry = 'point'
 par%spectral_type  = 'voigt'
 par%nxfreq         = 121
/
\end{lstlisting}

\subsection{Sphere with Hubble Outflow and Peel-off}

\begin{lstlisting}
&parameters
 par%no_photons     = 1e6
 par%temperature    = 1.0e4
 par%taumax         = 1.0e4
 par%rmax           = 1.0
 par%nx = 101, par%ny = 101, par%nz = 101
 par%source_geometry = 'point'
 par%spectral_type  = 'voigt'
 par%velocity_type  = 'hubble'
 par%Vexp           = 200.0          ! km/s outflow
 par%nxfreq         = 121
 par%nxim = 129, par%nyim = 129
 par%alpha = 0.0  45.0  90.0
 par%beta  = 0.0   0.0   0.0
 par%save_peeloff_2D = .true.
/
\end{lstlisting}

\subsection{Realistic Galaxy (FITS Density Field, Stokes Polarization)}

\begin{lstlisting}
&parameters
 par%no_photons    = 1e8
 par%dens_file     = 'galaxy.dens.fits.gz'
 par%temp_file     = 'galaxy.temp.fits.gz'
 par%velo_file     = 'galaxy.velo.fits.gz'
 par%taumax        = 1.0e4           ! rescale to this optical depth
 par%distance_unit = 'kpc'
 par%source_geometry = 'exponential'
 par%source_rscale = 3.0
 par%source_zscale = 0.5
 par%spectral_type = 'voigt'
 par%DGR           = 1.0
 par%use_stokes    = .true.
 par%nxfreq        = 121
 par%nxim = 200, par%nyim = 200
 par%save_peeloff_2D = .true.
/
\end{lstlisting}

\subsection{AMR Sphere (Generic Text Format)}

\begin{lstlisting}
&parameters
 par%use_amr_grid   = .true.         ! or par%grid_type = 'amr' in v2.10
 par%amr_type       = 'generic'
 par%amr_file       = 'sphere_amr.dat'
 par%no_photons     = 1e6
 par%taumax         = 1.0e4
 par%source_geometry = 'point'
 par%spectral_type  = 'voigt'
 par%nxfreq         = 121
 par%nxim = 100, par%nyim = 100
 par%save_peeloff_2D = .true.
 par%save_Jin       = .true.
/
\end{lstlisting}

\subsection{Clumpy Sphere with Random Velocities}

\begin{lstlisting}
&parameters
 par%use_clump_medium = .true.       ! or par%grid_type = 'clump' in v2.10
 par%rmax             = 1.0
 par%clump_radius     = 0.01
 par%clump_f_cov      = 5.0
 par%clump_tau0       = 1.0e3
 par%temperature      = 1.0e4
 par%clump_sigma_v    = 50.0         ! km/s random dispersion
 par%no_photons       = 1e6
 par%spectral_type    = 'voigt'
 par%nxfreq           = 301
 par%velocity_min = -200.0
 par%velocity_max =  200.0
 par%save_clump_info  = .true.
/
\end{lstlisting}

\end{document}
